\section{Conclusion}
\label{sec:conclusion}
We followed one edge node through learning, risk and explanation on C-MAPSS. For learning, a paired FD001--FD004 suite compared centralized, local-only and federated RUL prediction under whole-engine IID, window-level operating-regime and engine-count-imbalanced partitions. FedProx was close to FedAvg, the tested FedAdam configuration degraded, and the capped exponential score objective traded RMSE for lower mean NASA score; no paired test reached the 0.05 threshold with five seeds. For risk, quadratic \texttt{AsymmetricRULLoss} reduced late-prediction rates by \AsymLateReductionMin{}\%--\AsymLateReductionMax{}\% and decentralized NASA scores by \AsymDecentralNasaReductionMin{}\%--\AsymDecentralNasaReductionMax{}\% in a separate controlled FD001/FD002 comparison, while increasing symmetric RMSE on FD002. For explanation, a local attribution and drift module placed the core first in \DiagFdOneAttrCore{}\% and \DiagFdTwoAttrCore{}\% of single-fault low-RUL windows, and its signed drift profiles separated two seed-stable engine groups on FD003/FD004 with opposite sensor signatures. Without component labels, these are consistency and separability results, not diagnosis accuracy. The module made no socket calls and cost at most \DiagLatFullMax{}~ms per window on one thread of an Apple Silicon host. The committed edge artifacts come from Mac hosts, not Raspberry Pi hardware. Future work should measure both training and diagnostics on physical edge hardware, evaluate attribution against labeled component faults, and add authenticated transport and privacy-enhanced aggregation.
